\section{Results II: Whether the Faithful Geometry Is Used}
\label{sec:gap}
% ------------------------------------------------------------
% Sumber: finding 07 (weak instrument) + 08 (strong instrument, 3 tipe)
% + 09/11 (sweep 6 tipe) + 10v2 (probe). Judul lama "Causal Use Tracks
% Fidelity" DIGANTI: klaim itu gugur di 6 tipe (Sec. sweep).
% ------------------------------------------------------------

\subsection{Outputs are nearly identity-invariant under natural prompting}
\label{sec:ceiling}
% finding 07: full-sentence swap ceiling + D1-D2 — SUDAH FINAL
Before asking whether any internal location causally drives answers, we
measure the ceiling: how much does the model's predicted opinion
distribution move when the \emph{entire} demographic identity in the
prompt is replaced? Averaged over pairs and questions, the answer is
almost not at all: $\wdist(\hat{p}_A, p^{\mathrm{real}}_B) -
\wdist(\hat{p}_B, p^{\mathrm{real}}_B) \approx 0.004$ against a total
prediction error of ${\sim}0.37$ (${<}2\%$). (One illustrative item, not
an aggregate statistic: an 11-point real gap between 18--29- and
50--64-year-olds moves by only 1.7 points under a full persona
swap---${\sim}15\%$ of the true difference.)

Worse, the small variation that does exist is directionally
uninformative: the prediction for group $A$ is \emph{no closer} to $A$'s
own ground truth than to another group's (median advantage
$\approx0$; sign test ${\approx}$ coin flip in 2 of 3 tested types).
The model produces slightly-different-but-equally-wrong answers per
persona. % D2

We re-measured this ceiling on the instruction-tuned
Mistral-7B-Instruct-v0.2 under its own chat template: the near-invariance
replicates, with the instruct model's full-identity swap moving
predictions by $3.9\%$ of its total error (Appendix~\ref{app:ceiling-instruct};
Table~\ref{tab:ceiling}). This near-invariance is a property of the
model family, not of the raw-completion interface, though both estimates
assume the letter-probability read-out (Sec.~\ref{sec:limitations}).

The aggregate also hides a large per-type spread that inverts the
fidelity ranking at the two extremes (Appendix~\ref{app:ceiling-spread}).

\subsection{Patching at layer 11: a real but small effect}
\label{sec:weakinstrument}
% finding 07 (v1) + finding 08 (v2): FINAL. Ini subbab paling penting di Act 3.
Patching \starhead{} activations between group prompts (donor$\to$base) at
the \emph{prediction} token, with 1--3 heads at natural scale, moves
nothing (protocol in Appendix~\ref{app:patching}). But the ceiling above
shows this instrument had a dynamic range of under $2\%$ of total
prediction error to begin with---any causal effect this small would be
undetectable regardless of whether it exists.

\label{sec:v2}
We therefore repeat the intervention with an instrument matched in
strength to prior successful steering work \citep{llmopinions2026}: all
32 heads of layer 11 are patched jointly at a single identity token, so
the edit propagates to every later token (protocol in
Appendix~\ref{app:patching}). Across the three types tested here the
result is \textbf{graded}, and the gradation appears to line up with
each type's fidelity from Sec.~\ref{sec:act2}; Sec.~\ref{sec:sweep}
extends the experiment to all six types and shows that this alignment
does \emph{not} generalise.

\begin{table}[t]
\centering
\caption{Causal patching results with the strong instrument. Ceiling =
mean full-prompt-swap shift (Sec.~\ref{sec:ceiling}, this instrument's
question set). Shift = mean directional shift toward the target group's
true distribution; \% of ceiling in parentheses where meaningful.
$p_{\text{item}}$ = Wilcoxon vs.\ random-head control over 240 items
(descriptive; items are clustered); $p_{\text{pair}}$ = exact sign-flip
over the 12 pairs (primary).}
\label{tab:patchv2}
\resizebox{\linewidth}{!}{%
\begin{tabular}{lccccc}
\toprule
Type (split-half fidelity) & Ceiling & 32 heads @ L11 & Single \starhead{} & $p_{\text{item}}$ (32h) & $p_{\text{pair}}$ (32h) \\
\midrule
AGE$\times$POLPARTY ($0.63$)   & $+0.0165$ & $+0.0067$ (41\%) & $-0.0001$ (n.s.) & $0.0027$ & $0.019$ \\
RELIG$\times$POLPARTY ($0.52$) & $+0.0020$ & $+0.0034$ (170\%, n.s.) & $-0.0004$ (n.s.) & $0.16$ & $0.15$ \\
RACE$\times$RELIG ($0.21$)     & $+0.0026$ & $+0.0017$ (n.s.) & $\mathbf{-0.0005}$ (wrong dir., $p_{\text{pair}}{=}1/4096$) & $0.31$ & $0.29$ \\
\bottomrule
\end{tabular}%
}
\end{table}

For AGE$\times$POLPARTY, jointly patching all 32 heads at layer 11 shifts
predictions toward the target group's true distribution: $59\%$ of
trials move in the correct direction (vs.\ $32\%$ under an empirical
random-head control), recovering $41\%$ of the achievable ceiling. The
primary, pair-level test is marginal ($p_{\text{pair}}=0.019$; adding
layer 18: $0.016$) and does not survive a Holm correction over the
family of 24 tests behind Table~\ref{tab:patchv2}. Identity is thus only
\emph{plausibly} load-bearing here; Sec.~\ref{sec:sweep} localises
further.

For RACE$\times$RELIG, the 32-head patch is not significant
($p_{\text{pair}}=0.29$), while patching \starhead{} \emph{alone} shifts
predictions in the \emph{wrong} direction with the most consistent effect
in the table (worse than the random-head control in all 12 pairs;
$p_{\text{pair}}=1/4096$). A donor-control experiment shows this
backfire requires a coherent identity donor rather than merely a large
disruption to the representation (Appendix~\ref{app:amplification}).

Fidelity and single-unit causal sufficiency are dissociable:
\starhead{} alone never produces a correct, significant shift in any
type, and a correct effect requires the joint action of all 32
heads---faithful as a representation, but insufficient (or actively
misleading) as a sole intervention target. Scaling the patch further
backfires rather than strengthens it (Appendix~\ref{app:amplification}).

\subsection{The causal locus does not track fidelity}
\label{sec:sweep}
% Sumber: finding 09 (notebook 13) + finding 11 (notebook 14 Part B).
% Koreksi max-stat sign-flip 2000x, SATU run/seed utk keenam tipe:
% notebooks/output/14_.../sweep_all6_maxstat.csv + fixed_layer_L11_L1.csv
Section~\ref{sec:v2} localises an effect at layer 11 for one type and
finds none for another---too few types to tell whether causal use tracks
fidelity at all. We therefore replace the hypothesis-driven probe with
an exhaustive one, the same 32-head identity-token intervention applied
at \emph{every} layer, for all six attribute types (procedure and
selection corrections in Appendix~\ref{app:sweep-procedure}).

\begin{figure}[t]
\centering
\includegraphics[width=\linewidth]{fig_causal_sweep}
\caption{Causal effect of patching all 32 heads of each layer in turn, at
the identity token ($n{=}240$ pair--question items per type).
\emph{Descriptive} view: bars are the item-level paired $t$ of the
directional shift toward the target group's true distribution against a
random-layer control; the dashed line is the item-level max-statistic
noise ceiling. Inference in the text and Table~\ref{tab:sweep} is
cluster-robust (sign flips at the level of the 12 pairs).}
\label{fig:sweep}
\end{figure}

\begin{table}[t]
\centering
\caption{Per-layer patching sweep, all six types. Fidelity = split-half
median (one estimator used consistently; see text for sensitivity).
Item-level columns are descriptive ($t$ over 240 items; max-statistic $p$
over 32 layers); the primary inference is the pair-level max-statistic
$p$ (sign flips over 12 pairs, exact) and the pair-level $p$ at the
a-priori locus L11 (no layer selection). The item-level and pair-level
winning layers can differ, since averaging within pairs reweights the
evidence. Output ceiling = mean shift from replacing the entire
identity.}
\label{tab:sweep}
\resizebox{\linewidth}{!}{%
\begin{tabular}{lcccccccc}
\toprule
& & \multicolumn{3}{c}{item-level (descriptive)} & \multicolumn{3}{c}{pair-level (primary)} & \\
\cmidrule(lr){3-5}\cmidrule(lr){6-8}
Type & Fidelity & Win layer & $t$ & $p_{\max}$ & $t_{\text{win}}$ & $p_{\max}$ & $t$@L11 $\rightarrow$ $p$ & Out.\ ceiling \\
\midrule
RACE$\times$POLIDEOLOGY  & $0.30$ & L11 & $4.55$ & $0.0005$ & $3.22$ & $0.066$ & $\mathbf{3.22 \rightarrow 0.0020}$ & $0.0176$ \\
RACE$\times$RELIG        & $0.21$ & L1  & $4.10$ & $0.0005$ & $2.94$ (L8) & $0.122$ & $0.68 \rightarrow 0.27$ & $0.0026$ \\
AGE$\times$POLPARTY      & $0.63$ & L15 & $3.31$ & $0.0150$ & $2.37$ (L14) & $0.266$ & $1.93 \rightarrow 0.022$ & $0.0165$ \\
RELIG$\times$POLPARTY    & $0.52$ & L8  & $2.94$ & $0.0355$ & $2.58$ (L8) & $0.210$ & $1.11 \rightarrow 0.15$ & $0.0020$ \\
EDUCATION$\times$INCOME  & $0.58$ & L9  & $2.47$ & $0.1425$ & $1.57$ (L15) & $0.822$ & $0.75 \rightarrow 0.23$ & $0.0194$ \\
RACE$\times$POLPARTY     & $0.52$ & L14 & $2.40$ & $0.1770$ & $2.00$ (L4) & $0.524$ & $1.28 \rightarrow 0.11$ & $0.0032$ \\
\bottomrule
\end{tabular}%
}
\end{table}

Figure~\ref{fig:sweep} shows the descriptive item-level picture; the
pair-level statistics behind Table~\ref{tab:sweep} are conservative (12
clusters per type) and prune it sharply. No type survives the pair-level
max-statistic correction (best: RACE$\times$POLIDEOLOGY, $p=0.066$). The
a-priori test at L11 (no layer selection; fixed by Sec.~\ref{sec:act2})
is cleaner: RACE$\times$POLIDEOLOGY reaches $t=3.22$, exact $p=0.0020$,
the only entry that also survives a Bonferroni correction over the 12
fixed-location tests below (Table~\ref{tab:fixedlayer}); AGE$\times$POLPARTY
at L11 ($p=0.022$) is suggestive but not conclusive. The item-level loci
otherwise cluster mid-stack and at layer~1 (selection caveat in
Appendix~\ref{app:sweep-procedure}).

\begin{table}[t]
\centering
\caption{Fixed-location causal test, pair-level (exact sign-flip over 12
pairs, one-sided). L11 was fixed by Sec.~\ref{sec:act2} before any
sweep; L1 is pre-registered only for the upper block (see text).
Bonferroni for the 12 tests is $0.0042$; the single bold entry survives
it.}
\label{tab:fixedlayer}
\begin{tabular}{lcccc}
\toprule
Type & $t$ @ L11 & $p_{\text{pair}}$ & $t$ @ L1 & $p_{\text{pair}}$ \\
\midrule
\multicolumn{5}{l}{\emph{types not in the sweep that suggested L1:}} \\
RACE$\times$POLIDEOLOGY  & $\mathbf{3.22}$ & $\mathbf{0.0020}$ & $-0.74$ & $0.77$ \\
RACE$\times$POLPARTY     & $1.28$ & $0.114$ & $1.97$ & $0.033$ \\
EDUCATION$\times$INCOME  & $0.75$ & $0.234$ & $1.24$ & $0.123$ \\
\midrule
\multicolumn{5}{l}{\emph{types that were in that sweep (L1 column not selection-free):}} \\
AGE$\times$POLPARTY      & $1.93$ & $0.022$ & $0.34$ & $0.386$ \\
RELIG$\times$POLPARTY    & $1.11$ & $0.150$ & $1.46$ & $0.089$ \\
RACE$\times$RELIG        & $0.68$ & $0.267$ & $2.46$ & $0.016$ \\
\bottomrule
\end{tabular}
\end{table}

The natural hypothesis after Sec.~\ref{sec:v2}---that causal strength
should track fidelity---does not survive the extension to six types: the
rank correlation between fidelity and pair-level causal strength is
nowhere significantly positive under any estimator we have, including
the lexical-partial fidelity of Sec.~\ref{sec:lexical} ($+0.26$ to
$-0.77$ depending on ruler and layer, all $p \ge 0.07$; weakly powered at
$n{=}6$, Figure~\ref{fig:dissoc}). What carries the weight is the two
extreme types, tested \emph{directly} rather than by comparing
significance verdicts \citep{gelman2006}: the top-fidelity
types---whichever estimator names them---show no correction-surviving
single-layer locus (EDUCATION$\times$INCOME has the largest output
ceiling and largest raw shift of any type, yet $p=0.82$ at its best
layer; AGE$\times$POLPARTY, the split-half leader, reaches only a
suggestive $p=0.022$ at L11 that fails every correction here), while
RACE$\times$POLIDEOLOGY---among the weakest maps under every
estimator---is the single result that survives everything ($t=3.22$ at
L11, exact $p=0.0020$, surviving Bonferroni, recovering $42\%$ of its
ceiling). Representational fidelity is therefore not what decides
whether identity is causally used: the clearest causal pathway sits in
the type with the poorest map, and the graded pattern of Sec.~\ref{sec:v2}
does not generalise to all six types.

\begin{figure}[t]
\centering
\includegraphics[width=\linewidth]{fig_fidelity_vs_causal}
\caption{Representational fidelity (split-half median) against
cluster-robust causal strength (pair-level $t$), one point per attribute
type. Left: at the a-priori locus L11 (filled = exact pair-level
$p<0.05$). Right: at each type's own best layer (filled = survives the
pair-level max-statistic correction; none does).}
\label{fig:dissoc}
\end{figure}

Sec.~\ref{sec:act2}'s fidelity is scored on an identity-only prompt,
whereas the causal and probe experiments are scored in a full QA
context, so a tempting alternative reading is that the geometry simply
does not survive into the answering context, not that the model fails
to consult its map. Measured directly in that context---on the same
activations collected for the probe (Sec.~\ref{sec:probe})---the map
does weaken in five of six types,
but the dissociation itself persists under the QA-context ruler at the
same a-priori locus ($\rrho=+0.03$, $p=0.96$), so ``we measured the
wrong map'' does not explain the dissociation away
(Appendix~\ref{app:qa-context-map}).

\subsection{Reading the map directly: probe versus output}
\label{sec:probe}
% Sumber: finding 10 v2 (notebook 14 Part A + fitting lokal).
% Angka: notebooks/output/14_.../probe_v2_summary.csv, probe_v2_grouprank.csv,
% grouprank_noise_ceiling.csv
The causal experiments ask whether the model \emph{uses} what it
encodes; the complementary question is how much there is to use if we
read the faithful location directly rather than listen to the model's
mouth. We fit a ridge probe, evaluated leave-one-cell-out, at three
locations---full layer 11, the single head \starhead{}, and layer
1---against two references: the model's own answer, and a group-blind
baseline built from other cells' real answers on the same question
(full protocol in Appendices~\ref{app:probe-procedure}, \ref{app:probe}).

\begin{figure}[t]
\centering
\includegraphics[width=\linewidth]{fig_probe}
\caption{Left: distance to survey truth (Wasserstein, lower is better) for
the model's own answer, probes read from the faithful location, and a
group-blind baseline built from other cells' real answers. Right: how much
of the true per-question group \emph{ordering} each read-out recovers
(mean per-question Spearman across cells, half-split evaluation; error
bars are standard errors), against the reliability ceiling of the metric
given survey sampling noise.}
\label{fig:probe}
\end{figure}

In all six types the probe is substantially closer to survey truth than
the model's own answer ($22$--$30\%$ closer overall, better on
$65$--$68\%$ of individual items), not merely by better calibration:
fitting a temperature parameter to the mouth barely helps
($0.300\to0.298$ for AGE$\times$POLPARTY), reproducing against real
survey ground truth the claim that internal states are more informative
than generated answers \citep{llmopinions2026}.

But the gain is question-level, not group-level: a group-blind baseline
that ignores who is being asked and predicts the mean real distribution
of the other cells on the same question beats every read-out we have, by
a wide margin ($0.133$ vs.\ $0.211$ for AGE$\times$POLPARTY; $0.135$
vs.\ $0.234$ for EDUCATION$\times$INCOME). Ranking cells within each
question by predicted opinion and correlating with the true ranking
(half-split, to avoid a leave-one-out artifact) gives a flat null for
everything: the full-layer L11 probe recovers only $+0.02$ to $+0.07$,
the single head \starhead{} is noisier still ($-0.05$ to $+0.07$), and
the model's own answers recover $+0.06$ to $+0.11$---at least as much as
either probe---against a reliability ceiling of $+0.81$ to $+0.86$.

These two results measure different resolutions and are consistent:
Sec.~\ref{sec:act2} scores group geometry aggregated over hundreds of
questions, where the arrangement mirrors survey structure ($\rrho$ up to
$0.63$ held-out); this section scores per-item discrimination, where the
signal is near zero. The model holds a group map that is faithful on
average, but not sharp enough to say who answers what on any particular
question.

If the goal is a per-group opinion distribution, reading $128$
dimensions of one attention head is markedly better than asking the
model, but still worse than using other groups' real answers to the
same question (practical implications in Sec.~\ref{sec:discussion}).
